\section{Conclusion and Future Work}
\label{sec:conclusion}

In this paper, we introduced \textbf{PAC-Bayesian Smooth Trajectory Merging (PAC-STM)}, a rigorous learning-theoretic framework for dynamic layer-wise model merging in deep networks. By rejecting heuristic ensembling regularizers, we demonstrated that modeling layer-specific log-temperatures as a depth-wise Gaussian random walk prior allows the exact derivation of a closed-form Kullback-Leibler (KL) complexity penalty. This complexity penalty mathematically acts as a first-order ensembling smoothness regularizer, restricting the high-frequency parameter oscillations that cause transductive overfitting on ultra-low calibration sets ($N=16$).

Empirically, our results on a 14-layer Coordinate Sandbox demonstrate that PAC-STM successfully stabilizes deep ensembling trajectories, outperforming standard unregularized Empirical Risk Minimization (ERM) under heterogeneous batching workloads. Furthermore, by utilizing activation-blending, PAC-STM maintains absolute, provable immunity to both heterogeneity collapse and vectorization collapse under realistic edge-serving streaming workloads. 

Our work bridges the gap between randomized generalization bounds and stable ensembling trajectories. While the standard SVD-based UN-PCA-SEP projection assumes a linear subspace, we have successfully formulated, implemented, and empirically validated both the non-linear projection extension (UN-KPCA-SEP via Kernel PCA) and skip-connection-aware (residual) prior topologies. Our empirical evaluations demonstrate that both extensions yield substantial gains in ensembling accuracy and trajectory smoothness. In future work, we plan to focus on scaling PAC-STM to massive decoder-only large language models (LLMs) and multi-modal transformers where layer-wise adapter routing is increasingly critical. Additionally, we aim to explore more general graph-structured prior topologies, such as hierarchical trees or sparse graph Laplacians, to model ensembling trajectories over extremely large and complex task taxonomies.
